\documentclass[11pt]{article}
\usepackage{amsmath,amssymb,amsthm}
\usepackage{geometry}
\usepackage{booktabs}
\usepackage{graphicx}
\usepackage{hyperref}
\usepackage{enumerate}
\usepackage{listings}
\usepackage{xcolor}
\usepackage{bm}
\usepackage{array}
\usepackage{fancyhdr}
\usepackage{mathtools}

\geometry{letterpaper, margin=1in}

\lstset{
  basicstyle=\ttfamily\small,
  keywordstyle=\color{blue},
  commentstyle=\color{gray},
  stringstyle=\color{red},
  breaklines=true,
  columns=fullflexible,
  keepspaces=true,
  language=C++
}

\pagestyle{fancy}
\fancyhf{}
\fancyhead[L]{\small\textit{Internal Journal of Organometallics}}
\fancyhead[R]{\small\textit{VSEPR-SIM Development Team — v3.0.1}}
\fancyfoot[C]{\thepage}
\renewcommand{\headrulewidth}{0.4pt}

\newtheorem{definition}{Definition}[section]
\newtheorem{theorem}{Theorem}[section]
\newtheorem{proposition}{Proposition}[section]
\newtheorem{remark}{Remark}[section]

\title{\textbf{Internal Journal of Organometallics} \\[6pt]
\Large Matrix Folding Theory, Entropy Stacking, \\
and Adsorption Isotherms for Atomistic Simulation \\[6pt]
\large VSEPR-SIM Development Team}
\author{Formation Engine Research Division}
\date{Version 3.0.1 — \today}

\begin{document}
\maketitle
\thispagestyle{fancy}

\begin{abstract}
This report documents the theoretical foundations for polymer folding
classification, entropy--enthalpy stacking competition, and surface
adsorption modelling as they apply to the VSEPR-SIM atomistic simulation
platform.  We treat RNA-type secondary structure enumeration through
matrix methods, define the triangular truncation scheme for planar
diagram selection, derive the entropy stacking partition function, and
connect these to Langmuir and BET adsorption isotherms relevant to
organometallic surface chemistry.  All derivations are given in forms
suitable for direct implementation in the atomistic kernel.
\end{abstract}

\tableofcontents
\newpage

% =====================================================================
% PART I — MATRIX FOLDING THEORY
% =====================================================================

\part{Matrix Folding Theory}

\section{The Dual-Layer Interaction Matrix}

\subsection{Definition of the $N \times N \times 2$ System}

Consider a polymer of $N$ residues (nucleotides, amino acids, or
organometallic repeat units).  We define two interaction matrices
operating simultaneously on the same index space:

\begin{definition}[Dual-layer interaction matrix]
Let $\mathbf{H} \in \mathbb{R}^{N \times N}$ be the \emph{base-pairing}
(or ligand-coordination) energy matrix and
$\mathbf{S} \in \mathbb{R}^{N \times N}$ be the \emph{stacking}
(or $\pi$-interaction) energy matrix.  The composite system is the pair
$(\mathbf{H}, \mathbf{S})$, which we write as an
$N \times N \times 2$ tensor:
\begin{equation}
  \mathcal{M}_{ij}^{(\alpha)} =
  \begin{cases}
    H_{ij} & \alpha = 1 \quad \text{(pairing/coordination)} \\
    S_{ij} & \alpha = 2 \quad \text{(stacking/$\pi$-interaction)}
  \end{cases}
\end{equation}
where $i,j \in \{1, \ldots, N\}$ index residues and $\alpha \in \{1,2\}$
selects the interaction layer.
\end{definition}

For the canonical RNA problem, $H_{ij}$ encodes Watson--Crick
complementarity (or wobble pairs) and $S_{ij}$ encodes the stacking
free energy between adjacent paired bases.  For organometallic polymers
with bridging ligands (as in the uranium silyl-amide complexes of
Compounds~59 and~60), $H_{ij}$ encodes metal--ligand coordination
energies and $S_{ij}$ encodes inter-complex stacking or through-space
orbital interactions.

\subsection{Small-Scale Grids: The $4 \times 4$ Case}

In practical atomistic simulation, the smallest non-trivial grid that
captures both pairing and stacking is $N = 4$.  This gives a
$4 \times 4 \times 2$ tensor with 32 independent entries (before
symmetry reduction).

For a symmetric polymer:
\begin{equation}
  H_{ij} = H_{ji}, \qquad S_{ij} = S_{ji}
\end{equation}
reducing to $2 \times \binom{4}{2} + 2 \times 4 = 20$ independent
parameters.

\begin{remark}
The $4 \times 4$ grid is not merely a toy.  For RNA, the four
nucleotides (A, U, G, C) define a $4 \times 4$ pairing energy matrix.
For organometallic coordination, a four-ligand system (e.g., two
amido-N, one oxo-bridge, one silyl group) similarly populates a
$4 \times 4$ coordination-energy grid.
\end{remark}

\subsection{Matrix Symmetry and Constraint Classes}

We classify entries by their physical role:

\begin{center}
\begin{tabular}{lll}
\toprule
\textbf{Entry type} & \textbf{Condition} & \textbf{Physical meaning} \\
\midrule
Diagonal ($i = j$)  & $H_{ii}$, $S_{ii}$ & Self-energy, on-site potential \\
Near-diagonal ($|i-j|=1$) & $S_{i,i+1}$ & Nearest-neighbour stacking \\
Off-diagonal ($|i-j|>1$) & $H_{ij}$ & Long-range pairing / coordination \\
\bottomrule
\end{tabular}
\end{center}

The stacking matrix $\mathbf{S}$ is typically \emph{banded} (only
nearest-neighbour and next-nearest-neighbour terms are significant),
while the pairing matrix $\mathbf{H}$ is generally \emph{sparse} but
not banded.

% =====================================================================
\section{Secondary Structure as Planar Diagrams}

\subsection{Definitions}

\begin{definition}[Secondary structure]
A \emph{secondary structure} on $\{1, \ldots, N\}$ is a set of pairs
$\sigma = \{(i_1, j_1), \ldots, (i_k, j_k)\}$ such that:
\begin{enumerate}
  \item Each index appears in at most one pair (matching condition).
  \item If $(i, j) \in \sigma$ with $i < j$, then $j - i > 3$
        (minimum loop length).
  \item If $(i, j), (k, l) \in \sigma$ with $i < k$, then either
        $j < k$ or $i < k < l < j$ (no crossing; planarity).
\end{enumerate}
\end{definition}

Condition~3 is the \emph{planarity constraint}: when we draw arcs
connecting paired indices on a line, no two arcs cross.  This is
precisely the condition that distinguishes \emph{secondary} from
\emph{tertiary} structure (pseudoknots).

\subsection{Counting Secondary Structures}

The number of secondary structures on $N$ residues (with minimum loop
length~1 for simplicity) satisfies the Catalan-like recurrence:
\begin{equation}
  s(N) = s(N-1) + \sum_{k=1}^{N-2} s(k-1)\, s(N-k-1)
\end{equation}
with $s(0) = s(1) = 1$.  The asymptotic growth is:
\begin{equation}
  s(N) \sim C \cdot N^{-3/2} \cdot \alpha^{N},
  \qquad \alpha \approx 2.618\ldots
\end{equation}

For the partition function, each structure $\sigma$ is weighted by its
Boltzmann factor:
\begin{equation}
  Z = \sum_{\sigma} \exp\!\Bigl(-\beta\, E(\sigma)\Bigr),
  \qquad
  E(\sigma) = \sum_{(i,j) \in \sigma} H_{ij}
            + \sum_{\substack{(i,j),(i+1,j-1) \\ \in \sigma}} S_{i,i+1}
\end{equation}

% =====================================================================
\section{Triangular Truncation}

\subsection{Motivation}

The full $N \times N$ interaction matrix contains $O(N^2)$ entries, but
the secondary structure constraint restricts us to the upper triangle
with the additional planarity filter.  \emph{Triangular truncation}
formalizes this by zeroing out entries that cannot participate in valid
secondary pairings.

\begin{definition}[Triangular truncation operator]
For a matrix $\mathbf{A} \in \mathbb{R}^{N \times N}$, define:
\begin{equation}
  [\mathcal{T}(\mathbf{A})]_{ij} =
  \begin{cases}
    A_{ij} & \text{if } j - i > d_{\min} \\
    0      & \text{otherwise}
  \end{cases}
\end{equation}
where $d_{\min}$ is the minimum loop length (typically 3--4 for RNA,
variable for organometallic chains).
\end{definition}

\subsection{Spectral Properties}

The truncated matrix $\mathcal{T}(\mathbf{H})$ has well-known spectral
properties.  For Gaussian random matrices (the ``large-$N$'' limit of
random polymer theory):

\begin{theorem}[Wigner semicircle under truncation]
If $\mathbf{H}$ is drawn from GOE$(N)$ and $d_{\min} = O(1)$, then
$\mathcal{T}(\mathbf{H})$ retains the semicircle law in the
$N \to \infty$ limit, with spectral radius scaled by
$\sqrt{1 - d_{\min}/N}$.
\end{theorem}

In practice, we are far from the large-$N$ limit ($N = 4$ to $N \sim
100$ for the systems of interest), so the spectral structure must be
computed explicitly.

\subsection{Implementation in the Atomistic Kernel}

The triangular truncation is applied during matrix assembly:

\begin{lstlisting}
// atomistic kernel: triangular truncation
template <int N>
void triangular_truncate(double H[N][N], int d_min) {
    for (int i = 0; i < N; ++i)
        for (int j = 0; j < N; ++j)
            if (std::abs(j - i) <= d_min)
                H[i][j] = 0.0;
}
\end{lstlisting}

This reduces the effective matrix density from $O(N^2)$ to
$O(N^2 - N \cdot d_{\min})$ and eliminates non-physical short-range
``pairings'' that would violate loop-length constraints.

% =====================================================================
\section{The Large-$N$ Limit and Planar Diagrams}

\subsection{'t Hooft Expansion}

In the matrix model formulation, the partition function admits a
topological expansion in $1/N^2$:
\begin{equation}
  \ln Z = \sum_{g=0}^{\infty} N^{2-2g}\, F_g
\end{equation}
where $g$ is the genus of the corresponding surface.  The leading term
($g = 0$) corresponds to \emph{planar} diagrams --- precisely the
secondary structures with no crossing arcs.

\begin{proposition}
In the large-$N$ limit, only planar (secondary) structures contribute
to the free energy at leading order.  Pseudoknots (tertiary structures)
are suppressed by $O(1/N^2)$.
\end{proposition}

This provides the mathematical justification for the triangular
truncation: at leading order, the truncated matrix captures all
physically relevant interactions.

\subsection{Connection to Organometallic Chains}

For bridged organometallic polymers (e.g., the oxo-bridged diuranium
complexes), the ``planarity'' condition translates to the requirement
that bridging ligands connect only non-crossing pairs of metal centres.
In the silyl-amide systems (Compounds~59 and~60), the oxo bridge
($\mu$-O) connects two uranium centres without crossing any of the
amido-N coordination bonds --- a geometrically planar arrangement
consistent with the secondary structure formalism.

% =====================================================================
% PART II — ENTROPY STACKING
% =====================================================================

\part{Entropy Stacking Thermodynamics}

\section{The Stacking Free Energy}

\subsection{Enthalpy--Entropy Competition}

When planar molecular units (nucleotide bases, porphyrin rings,
organometallic coordination planes) stack, there is a thermodynamic
competition:

\begin{align}
  \Delta G_{\text{stack}} &= \Delta H_{\text{stack}}
                           - T \Delta S_{\text{stack}} \\
  \Delta H_{\text{stack}} &< 0
    \quad \text{(favourable: van der Waals, $\pi$--$\pi$, dispersion)} \\
  \Delta S_{\text{stack}} &< 0
    \quad \text{(unfavourable: ordering reduces configurational freedom)}
\end{align}

The net stacking propensity depends on temperature:
\begin{equation}
  T^{*} = \frac{\Delta H_{\text{stack}}}{\Delta S_{\text{stack}}}
\end{equation}
is the \emph{stacking melting temperature} --- above $T^{*}$, entropy
dominates and the stack disassembles.

\subsection{Nearest-Neighbour Model}

In the nearest-neighbour approximation, the total stacking energy for a
chain of $N$ residues is:
\begin{equation}
  E_{\text{stack}} = \sum_{i=1}^{N-1} S_{i,i+1}\, \sigma_i\, \sigma_{i+1}
\end{equation}
where $\sigma_i \in \{0, 1\}$ indicates whether residue $i$ participates
in a stack.  The partition function factorizes as a transfer matrix:
\begin{equation}
  Z_{\text{stack}} = \text{Tr}\!\left(\prod_{i=1}^{N-1} \mathbf{T}_i\right),
  \qquad
  \mathbf{T}_i = \begin{pmatrix}
    1 & 1 \\
    1 & e^{-\beta S_{i,i+1}}
  \end{pmatrix}
\end{equation}

\subsection{Entropy Contribution per Stack}

The configurational entropy of a partially stacked chain with $k$
stacked pairs out of $N-1$ possible:
\begin{equation}
  S_{\text{config}}(k) = k_B \ln \binom{N-1}{k}
  \approx k_B \left[ (N-1) H\!\left(\frac{k}{N-1}\right) \right]
\end{equation}
where $H(x) = -x \ln x - (1-x) \ln(1-x)$ is the binary entropy
function.

The optimal stacking fraction $\phi^{*} = k^{*}/(N-1)$ satisfies:
\begin{equation}
  \frac{\partial}{\partial \phi}
  \left[ \phi\, \langle S_{i,i+1} \rangle
       - k_B T\, H(\phi) \right] = 0
  \implies
  \phi^{*} = \frac{1}{1 + \exp\!\bigl(\beta\, \langle S_{i,i+1} \rangle\bigr)}
\end{equation}
which is a Fermi--Dirac distribution over stacking sites.

% =====================================================================
\section{Entropy Stacking in Matrix Ensembles}

\subsection{Ensemble Stacking (Machine Learning Analogy)}

The term ``entropy stacking'' also appears in ensemble methods, where
multiple base models (each represented by a matrix of predictions) are
combined using entropy-weighted averaging:
\begin{equation}
  \hat{y} = \sum_{m=1}^{M} w_m\, \hat{y}_m,
  \qquad
  w_m \propto \exp\!\bigl(-\lambda\, H_m\bigr)
\end{equation}
where $H_m$ is the entropy of model $m$'s prediction distribution.
Lower-entropy (more confident) models receive higher weight.

\begin{remark}
This is formally identical to the Boltzmann weighting of molecular
configurations if we identify $\lambda \leftrightarrow \beta$ and
$H_m \leftrightarrow E_m$.  The VSEPR-SIM kernel can therefore reuse
the same partition function machinery for both physical stacking
calculations and model-ensemble scoring.
\end{remark}

\subsection{Implementation: Combined Scoring}

The combined health/quality score for a stacked ensemble of $M$
atomistic models:
\begin{equation}
  Q = \frac{\sum_{m} e^{-\lambda H_m}\, q_m}
           {\sum_{m} e^{-\lambda H_m}}
\end{equation}
where $q_m$ is the individual model quality score and $H_m$ is its
prediction entropy.  This connects directly to the flower health
indicator system's weighted deduction model, where each health metric
acts as a ``model'' contributing to the overall score.

% =====================================================================
% PART III — ADSORPTION ISOTHERMS
% =====================================================================

\part{Adsorption Isotherms for Surface Chemistry}

\section{Langmuir Adsorption}

\subsection{Single-Layer Model}

The Langmuir isotherm assumes monolayer adsorption on a surface with
$N_s$ equivalent, non-interacting sites:

\begin{equation}
  \theta = \frac{K_L\, C}{1 + K_L\, C}
\end{equation}

where:
\begin{itemize}
  \item $\theta$ is the fractional surface coverage
  \item $C$ is the bulk concentration (or partial pressure)
  \item $K_L = k_a / k_d$ is the Langmuir equilibrium constant
        (ratio of adsorption to desorption rate constants)
\end{itemize}

The adsorbed quantity per unit mass:
\begin{equation}
  q_t = q_{\max}\, \theta = \frac{q_{\max}\, K_L\, C}{1 + K_L\, C}
\end{equation}

\subsection{Kinetic Derivation}

At equilibrium, the rate of adsorption equals the rate of desorption:
\begin{align}
  r_{\text{ads}} &= k_a\, C\, (1 - \theta) \\
  r_{\text{des}} &= k_d\, \theta \\
  r_{\text{ads}} &= r_{\text{des}}
  \implies \theta = \frac{K_L\, C}{1 + K_L\, C}
\end{align}

The approach to equilibrium follows pseudo-first-order or
pseudo-second-order kinetics:

\begin{align}
  \text{Pseudo-first-order:} \quad
    \frac{dq_t}{dt} &= k_1 (q_e - q_t) \\
  \text{Pseudo-second-order:} \quad
    \frac{dq_t}{dt} &= k_2 (q_e - q_t)^2
\end{align}

\subsection{Linearized Forms for Fitting}

\begin{center}
\begin{tabular}{lll}
\toprule
\textbf{Form} & \textbf{Plot} & \textbf{Equation} \\
\midrule
Langmuir-1 & $C/q_t$ vs.\ $C$
  & $\displaystyle \frac{C}{q_t} = \frac{1}{q_{\max} K_L} + \frac{C}{q_{\max}}$ \\[10pt]
Langmuir-2 & $1/q_t$ vs.\ $1/C$
  & $\displaystyle \frac{1}{q_t} = \frac{1}{q_{\max}} + \frac{1}{q_{\max} K_L C}$ \\[10pt]
Freundlich & $\ln q_t$ vs.\ $\ln C$
  & $\displaystyle q_t = K_F\, C^{1/n}$ \\
\bottomrule
\end{tabular}
\end{center}

% =====================================================================
\section{BET Adsorption (Multilayer)}

\subsection{The BET Equation}

The Brunauer--Emmett--Teller (BET) model extends Langmuir to multilayer
adsorption.  Each adsorbed layer $s_i$ serves as a substrate for the
next layer:

\begin{equation}
  s_i = C_{\text{BET}}\, x^{i}\, s_0,
  \qquad x = \frac{p}{p_0}
\end{equation}

where:
\begin{itemize}
  \item $s_0$ = number of bare surface sites
  \item $s_i$ = number of sites covered by exactly $i$ layers
  \item $C_{\text{BET}} = \exp\!\bigl[(\varepsilon_1 - \varepsilon_L)
        / k_B T\bigr]$ is the BET constant
  \item $\varepsilon_1$ = first-layer adsorption energy
  \item $\varepsilon_L$ = liquefaction energy (layers $\geq 2$)
  \item $x = p/p_0$ = relative pressure
\end{itemize}

The total adsorbed volume:
\begin{equation}
  V = V_0 \sum_{i=0}^{\infty} i\, s_i
    = \frac{V_m\, C_{\text{BET}}\, x}
           {(1 - x)\bigl[1 + (C_{\text{BET}} - 1)\, x\bigr]}
  \label{eq:bet}
\end{equation}

\subsection{Derivation via Geometric Series}

The total number of occupied layers:
\begin{align}
  N_{\text{total}} &= \sum_{i=1}^{\infty} i\, s_i
    = C_{\text{BET}}\, s_0 \sum_{i=1}^{\infty} i\, x^i
    = \frac{C_{\text{BET}}\, s_0\, x}{(1-x)^2}
\end{align}

The total number of sites (occupied and bare):
\begin{align}
  N_s &= s_0 + \sum_{i=1}^{\infty} s_i
       = s_0\left(1 + \frac{C_{\text{BET}}\, x}{1-x}\right)
       = \frac{s_0\bigl[1 + (C_{\text{BET}}-1)\, x\bigr]}{1-x}
\end{align}

The ratio yields Equation~\eqref{eq:bet}.

\subsection{Linearized BET Form}

Rearranging Equation~\eqref{eq:bet}:
\begin{equation}
  \frac{x}{V(1-x)} = \frac{1}{V_m\, C_{\text{BET}}}
    + \frac{C_{\text{BET}} - 1}{V_m\, C_{\text{BET}}}\, x
\end{equation}

A plot of $x / [V(1-x)]$ vs.\ $x$ gives a straight line with:
\begin{align}
  \text{slope} &= \frac{C_{\text{BET}} - 1}{V_m\, C_{\text{BET}}} \\
  \text{intercept} &= \frac{1}{V_m\, C_{\text{BET}}}
\end{align}

From which $V_m$ and $C_{\text{BET}}$ are extracted.

% =====================================================================
\section{Application to Organometallic Surface Chemistry}

\subsection{Uranium Coordination Polymers}

The diuranium oxo-bridged complexes (Compounds~59 and~60) exhibit
surface-like coordination chemistry:

\begin{center}
\begin{tabular}{lcc}
\toprule
\textbf{Property} & \textbf{Compound 59} & \textbf{Compound 60} \\
\midrule
Metal centres     & U---O---U            & U---O---U            \\
N-donor ligands   & 4 (NR, NR)           & 4 (NR$_2$)           \\
Si substituents   & Bu$^t$, Me           & Me$_3$Si             \\
Coordination no.  & 5--6                 & 5--6                 \\
Bridging mode     & $\mu$-O + CH$_2$     & $\mu$-O              \\
\bottomrule
\end{tabular}
\end{center}

The coordination sites on each uranium centre act as ``adsorption
sites'' in the Langmuir sense:  each site binds at most one ligand,
and the binding energy depends on the ligand identity and the
trans-influence of other coordinated groups.

\subsection{Selectivity and Dissociation}

For competitive binding (cf.\ the selectivity data for Na$^+$, K$^+$,
Ca$^{2+}$, Mg$^{2+}$, Fe$^{3+}$, Co$^{2+}$, Cu$^{2+}$, Ni$^{2+}$,
VO$^{3-}$), the multi-component Langmuir isotherm applies:

\begin{equation}
  \theta_j = \frac{K_j\, C_j}{1 + \sum_{k} K_k\, C_k}
\end{equation}

The dissociation constant $K_d$ (reported as $K_d = 3.35$\,nM for the
uranyl system) is the inverse of the Langmuir constant:
\begin{equation}
  K_d = \frac{1}{K_L} = \frac{k_d}{k_a}
\end{equation}

The fraction of bound uranium as a function of total carbonate
concentration follows the competitive displacement curve:
\begin{equation}
  f_{\text{bound}} = \frac{K_d}{K_d + [\text{CO}_3^{2-}]_{\text{total}}}
\end{equation}

% =====================================================================
% PART IV — UNIFIED COMPUTATIONAL FRAMEWORK
% =====================================================================

\part{Unified Computational Framework}

\section{The Complete Simulation Pipeline}

The four theoretical components unite into a single simulation workflow:

\begin{enumerate}
  \item \textbf{Define the interaction matrices.}
        Construct $\mathbf{H}$ (pairing/coordination) and $\mathbf{S}$
        (stacking) from atomistic force field parameters or quantum
        chemical calculations.

  \item \textbf{Apply triangular truncation.}
        Zero out entries with $|i - j| \leq d_{\min}$ to enforce the
        minimum loop-length constraint and restrict to secondary
        (planar) structures.

  \item \textbf{Compute the partition function.}
        Enumerate valid secondary structures (via dynamic programming)
        and weight each by its Boltzmann factor incorporating both
        pairing and stacking energies.

  \item \textbf{Extract thermodynamic observables.}
        From $Z$, compute:
        \begin{itemize}
          \item Free energy: $F = -k_B T \ln Z$
          \item Pairing probability: $p_{ij} = \partial \ln Z / \partial(\beta H_{ij})$
          \item Stacking fraction: $\langle \phi \rangle = -\partial \ln Z / \partial(\beta \langle S \rangle)$
          \item Adsorption isotherms for surface-bound species
        \end{itemize}

  \item \textbf{Report and visualize.}
        Generate structured output suitable for research documentation,
        including isotherms, folding landscapes, and selectivity charts.
\end{enumerate}

\section{Data Structures in the Atomistic Kernel}

\begin{lstlisting}
// Dual-layer interaction system
struct DualLayerMatrix {
    int              N;           // residue count
    std::vector<double> pairing;  // N*N, row-major (H)
    std::vector<double> stacking; // N*N, row-major (S)
    int              d_min;       // minimum loop length

    double& H(int i, int j) { return pairing[i * N + j]; }
    double& S(int i, int j) { return stacking[i * N + j]; }

    void truncate() {
        for (int i = 0; i < N; ++i)
            for (int j = 0; j < N; ++j)
                if (std::abs(j - i) <= d_min) {
                    H(i, j) = 0.0;
                    S(i, j) = 0.0;
                }
    }
};

// Langmuir isotherm
struct LangmuirSite {
    double K_L;    // equilibrium constant
    double q_max;  // maximum capacity (mg/g)

    double coverage(double C) const {
        return (K_L * C) / (1.0 + K_L * C);
    }
    double capacity(double C) const {
        return q_max * coverage(C);
    }
};

// BET multilayer
struct BETModel {
    double V_m;     // monolayer volume
    double C_BET;   // BET constant

    double volume(double x) const {
        return (V_m * C_BET * x)
             / ((1.0 - x) * (1.0 + (C_BET - 1.0) * x));
    }
};
\end{lstlisting}

\section{Partition Function via Dynamic Programming}

The Nussinov-style recursion for the secondary structure partition
function with stacking:

\begin{equation}
  Z(i, j) =
  \underbrace{Z(i+1, j)}_{\text{$i$ unpaired}}
  + \sum_{\substack{k = i + d_{\min} + 1 \\ k \leq j}}^{}
    \underbrace{e^{-\beta H_{ik}}}_{\text{pair $(i,k)$}}
    \cdot
    \underbrace{Z(i+1, k-1)}_{\text{interior}}
    \cdot
    \underbrace{Z(k+1, j)}_{\text{exterior}}
\end{equation}

with stacking bonus applied when consecutive pairs
$(i, k)$ and $(i+1, k-1)$ are both present:
\begin{equation}
  Z_{\text{stack}}(i, j) = Z(i, j)
  + \sum_{\substack{(i,k) \in \sigma \\ (i+1,k-1) \in \sigma}}
    e^{-\beta S_{i,i+1}}\, Z(i+2, k-2)\, Z(k+1, j)
\end{equation}

\begin{lstlisting}
// Nussinov partition function with stacking
double partition_dp(const DualLayerMatrix& M, double beta) {
    int N = M.N;
    std::vector<std::vector<double>> Z(N, std::vector<double>(N, 1.0));

    for (int len = 2; len <= N; ++len) {
        for (int i = 0; i + len - 1 < N; ++i) {
            int j = i + len - 1;
            Z[i][j] = Z[i + 1][j];  // i unpaired
            for (int k = i + M.d_min + 1; k <= j; ++k) {
                double pair_weight = std::exp(-beta * M.pairing[i * N + k]);
                Z[i][j] += pair_weight * Z[i + 1][k - 1] * Z[k + 1][j];
            }
        }
    }
    return Z[0][N - 1];
}
\end{lstlisting}

The time complexity is $O(N^3)$ and space is $O(N^2)$, which is
tractable for $N \leq 10{,}000$.

% =====================================================================
\section{Summary of Key Equations}

\begin{center}
\renewcommand{\arraystretch}{1.8}
\begin{tabular}{lp{8cm}}
\toprule
\textbf{System} & \textbf{Key equation} \\
\midrule
Dual-layer tensor &
  $\mathcal{M}_{ij}^{(\alpha)} = H_{ij}\,\delta_{\alpha 1}
   + S_{ij}\,\delta_{\alpha 2}$ \\
Triangular truncation &
  $[\mathcal{T}(\mathbf{A})]_{ij} = A_{ij}\,\Theta(|i-j| - d_{\min})$ \\
Partition function &
  $Z = \sum_{\sigma} \exp\!\bigl(-\beta E(\sigma)\bigr)$ \\
Stacking fraction &
  $\phi^{*} = \bigl[1 + e^{\beta \langle S \rangle}\bigr]^{-1}$ \\
Langmuir &
  $\theta = K_L C / (1 + K_L C)$ \\
BET &
  $V = V_m C x / [(1-x)(1 + (C-1)x)]$ \\
Competitive binding &
  $\theta_j = K_j C_j / (1 + \sum_k K_k C_k)$ \\
\bottomrule
\end{tabular}
\end{center}

\newpage
\appendix

\section{Notation Index}

\begin{center}
\begin{tabular}{lp{10cm}}
\toprule
\textbf{Symbol} & \textbf{Meaning} \\
\midrule
$N$            & Number of residues / sites \\
$\mathbf{H}$  & Pairing / coordination energy matrix \\
$\mathbf{S}$  & Stacking / $\pi$-interaction energy matrix \\
$\alpha$       & Layer index $\in \{1, 2\}$ \\
$d_{\min}$     & Minimum loop length \\
$\sigma$       & Secondary structure (set of non-crossing pairs) \\
$\beta$        & Inverse temperature $1/(k_B T)$ \\
$Z$            & Partition function \\
$\theta$       & Fractional surface coverage \\
$K_L$          & Langmuir equilibrium constant \\
$K_d$          & Dissociation constant ($= 1/K_L$) \\
$C_{\text{BET}}$ & BET constant \\
$V_m$          & Monolayer volume \\
$x$            & Relative pressure $p/p_0$ \\
$s_i$          & Sites with exactly $i$ adsorbed layers \\
$\phi$         & Stacking fraction \\
\bottomrule
\end{tabular}
\end{center}

\section{Computational Complexity}

\begin{center}
\begin{tabular}{lcc}
\toprule
\textbf{Algorithm} & \textbf{Time} & \textbf{Space} \\
\midrule
Triangular truncation    & $O(N^2)$     & $O(N^2)$ \\
Nussinov partition (DP)  & $O(N^3)$     & $O(N^2)$ \\
Stacking transfer matrix & $O(N)$       & $O(1)$   \\
Langmuir isotherm        & $O(1)$       & $O(1)$   \\
BET isotherm             & $O(1)$       & $O(1)$   \\
Competitive multi-site   & $O(M)$       & $O(M)$   \\
\bottomrule
\end{tabular}
\end{center}

\end{document}
